\documentclass[10pt,a4paper]{amsart}
\usepackage[margin=19mm]{geometry}
\usepackage{amsmath,amssymb,mathtools}
\usepackage[scr=rsfs,cal=euler]{mathalfa}
\usepackage{xfrac}
\usepackage[unicode,hidelinks,bookmarks=false]{hyperref}
\newcommand{\ie}{i.e.\ }
\newcommand{\cf}{cf.\ }
\newcommand{\resp}{respectively\ }
\newcommand{\Subsection}{\subsection}
\providecommand{\nobreakdash}{\nobreak\hbox{-}}
\setlength{\emergencystretch}{2em}
\multlinegap0pt
\usepackage{bm}
\usepackage{bbm}
\usepackage{dsfont}
\usepackage{xspace}
\usepackage{cancel}
\usepackage{mathtools}
\usepackage{amsthm}
\makeatletter
\@ifundefined{theorem}{\newtheorem{theorem}{Theorem}}{}
\@ifundefined{lemma}{\newtheorem{lemma}[theorem]{Lemma}}{}
\makeatother
\title[A New Spectral Conjugate Subgradient Method with Application]{A New Spectral Conjugate Subgradient Method with Application in Computed Tomography Image Reconstruction}
\author{Milagros Loreto, Thomas Humphries, Chella Raghavan, Kenneth Wu, Sam Kwak}
\date{Source: arXiv:2309.15266v2 (CC BY 4.0)}
\begin{document}
\maketitle

\noindent\emph{Source and licence.} A New Spectral Conjugate Subgradient Method with Application in Computed Tomography Image Reconstruction. Version arXiv:2309.15266v2 (CC BY 4.0), distributed under the Creative Commons Attribution 4.0 International licence, as recorded in the arXiv licence metadata for this exact version and shown on the abstract page https://arxiv.org/abs/2309.15266v2. Full rights notice: licenses/arxiv-2309.15266-CC-BY-4.0.txt.\par\smallskip

\noindent\textit{Editorial scope.} Counted unit: the Lemma whose source label is \texttt{L:lemma3} in the article (source0.tex lines 580--584, source label \texttt{L:lemma3}), delivered together with the article's own complete original proof of it (source0.tex lines 585--623, one \texttt{proof} environment of 39 lines). No mathematics is written, repaired or re-derived: statement and proof are the article's own text line for line. Required context retained uncounted: L:lemma2 (lines 552--557). Every definition, display and label of those lines is retained unchanged. Omitted from this package and not redistributed: 1--551, 558--579, 624--679. Nothing inside a retained excerpt is omitted or altered: every retained body is delivered exactly as the article's lines are written, so no omission receipt of any kind accompanies this package. Citations: the retained excerpts carry no citation command at all, so no bibliography entry of the article is transcribed and no reference is redistributed. Reference adaptation: none was needed and none was made; every reference inside the delivered unit resolves inside it, so no unresolved reference or citation is produced. Printed slips kept literal: no printed word, symbol, hypothesis or number of the counted statement or of its proof was altered. Layout and class: the article's own class is \texttt{interact} with packages including epstopdf, subfigure, xcolor, soul, natbib, url; the wrapper is the lane's pinned \texttt{amsart} editorial wrapper at 10pt on A4, which restates the theorem-like environments the retained text needs and adds the article's own macro definitions that the retained text needs (none), with extra wrapper packages (bm, bbm, dsfont, xspace, cancel, mathtools). The delivered numbering is the unit's own. Assets: no retained span contains a figure environment, a graphics-inclusion command, a PDF-page inclusion, a table or a TikZ picture; no image file is copied into this package, the package declares no asset dependency and no image file is redistributed with it. Licence: version arXiv:2309.15266v2 of this article is distributed under the Creative Commons Attribution 4.0 International licence, as recorded in the arXiv licence metadata for this exact version and shown on the abstract page https://arxiv.org/abs/2309.15266v2. A full rights notice is delivered at licenses/arxiv-2309.15266-CC-BY-4.0.txt. Characters: every retained excerpt is pure ASCII, so the delivered engine, XeLaTeX with its default Latin Modern text font, reports no missing character.
\begin{lemma}\label{L:lemma2}
Suppose $f: \mathbb{R}^n \to \mathbb{R}$ is subdifferentiable and its subgradient is Lipschitz continuous with constant $L$, and let $g \in \partial f(x)$. Then for all $x,p \in \mathbb{R}^n$,
\begin{align*}
\vert f(x+p) - f(x) - g^T p \vert \leq \frac{L}{2} \Vert p \Vert^2
\end{align*}
\end{lemma}

\begin{lemma}\label{L:lemma3} Suppose $f: \mathbb{R}^n \to \mathbb{R}$ is subdifferentiable and its subgradient is Lipschitz continuous with constant $L$. Under the nonmonotone line search condition, either $\alpha_k~=~1$, or we have the following lower bound on $\alpha_k$:
\begin{align*}
\alpha_k &\geq \frac{\gamma-1}{L} \frac{g_k^T d_k}{\Vert d_k \Vert^2}
\end{align*}
\end{lemma}

\begin{proof} Note that $\alpha_k = 1$ when $h_k = 0$, so we consider only the case $h_k \geq 1$, meaning $\alpha_k < 1$. There are two possibilities:
\begin{enumerate}
\item $\displaystyle \max_{0 \leq j \leq M} f(x_{k-j}) = f(x_k)$

In this case, since $\alpha_k < 1$, the line search condition must have been violated with step length $2 \alpha_k$:
\begin{align}
f(x_k + 2 \alpha_k d_k) - f(x_{k}) > 2\gamma \alpha_k g_k^T d_k + \eta_k.\label{E:l3e1}
\end{align}

By Lemma~\ref{L:lemma2}, we also have:
\begin{align}
f(x_k + 2 \alpha_k d_k) - f(x_{k}) -2  \alpha_k g_k^T d_k &\leq \frac{1}{2} L \Vert 2 \alpha_k d_k \Vert^2 \notag\\
f(x_{k}) - f(x_k + 2 \alpha_k d_k)  &\geq -2  \alpha_k g_k^T d_k - 2 \alpha_k^2 L \Vert d_k \Vert^2 -\eta_k,\label{E:l3e2}
\end{align}
where the absolute value can be removed on the left side since the quantity is positive, and $\eta_k$ can be subtracted on the second line because it is also positive. Adding (\ref{E:l3e1}) and (\ref{E:l3e2}) then gives:
\begin{align*}
0 &\geq 2\alpha_k (\gamma - 1)g_k^T d_k - 2 \alpha_k^2 L\Vert d_k \Vert ^2   \\
\alpha_k L\Vert d_k \Vert ^2 &\geq (\gamma - 1) g_k^T d_k \\
\alpha_k &\geq \frac{\gamma-1}{L} \frac{g_k^T d_k}{\Vert d_k \Vert^2},
\end{align*}
giving the desired lower bound. We now consider the second possibility:
\item  $\displaystyle \max_{0 \leq j \leq M} f(x_{k-j}) = f(x_{k-J})$, $1 \leq J \leq M$.

In this case, $\alpha_k$ is the first step size satisfying 
\begin{align*}
f(x_k + \alpha_k d_k ) \leq  f(x_{k-J}) + \gamma \alpha_k g_k^T d_k + \eta_k,
\end{align*}

Define $\alpha_k^*$ to be the step size that would be chosen in the case  $\displaystyle \max_{0 \leq j \leq M} f(x_{k-j}) = f(x_k)$, which was analyzed above. Suppose $\alpha_k < \alpha_k^*$, or equivalently, $\alpha_k^* = 2^n \alpha_k$ for some integer $n \geq 1$. Then it must be the case that
\begin{align*}
f(x_k + \alpha_k^* d_k) &> f(x_{k-J}) + \gamma \alpha_k^* g_k^T d_k + \eta_k
\end{align*}
since otherwise $\alpha_k^*$ would have been accepted as the step length, as it precedes $\alpha_k$ in the sequence $\{2^{-h_k}$\}. But since $f(x_{k-J}) \geq f(x_k)$, this implies
\begin{align*}
f(x_k + \alpha_k^* d_k) &> f(x_{k}) + \gamma \alpha_k^* g_k^T d_k + \eta_k,
\end{align*}
contradicting the definition of $\alpha_k^*$. Thus in this second case we must have $\displaystyle \alpha_k \geq \alpha^* \geq \frac{\gamma-1}{L} \frac{g_k^T d_k}{\Vert d_k \Vert^2}$ as well. Intuitively, the step size chosen by a nonmonotone line search with $M \geq 1$ can never be smaller than that chosen by the monotone (Armijo) line search corresponding to $M=0$, since the former is a less restrictive condition.
\end{enumerate}
\end{proof}

\end{document}
